\scanpage{25}
\diagram{\begin{tikzcd}[ampersand replacement=\&,column sep=normal,row sep=small]
\operatorname{Ker}i\arrow[d,very thick]\arrow[r,dashed] \& L\arrow[dl,dashed] \&\\
J \& L'\arrow[l,dashed]\arrow[u,dashed,"\unsure{2|\alpha|?}"'] \& M'\arrow[l,dashed]\\
\& \& \operatorname{Ker}\beta\arrow[u,dashed]\arrow[ull,thick]
\end{tikzcd}}

\diagram{\begin{tikzcd}[ampersand replacement=\&,row sep=small]
L+M'\arrow[r,very thick,"{(i,\beta)}","{\text{surj}}"' ]\arrow[d,"{(\sigma,\sigma\alpha\pi')}"',"{=\rho}"] \& M\arrow[ddl,dashed,bend left=12,"\mu"]\\
J\arrow[d,"\lambda"'] \&\\
G \&
\end{tikzcd}}

\margin{condition de non ramification}

\margin{le foncteur en $G$ est repr. par $\operatorname{Coker}\rho$, dont la formation commute à tout \unsure{chgt|chgt.} de base}

En résumé, la condition en question est l'annulation de $\lambda$ sur $R=R_1+R_2$, donc la factorisation de $\lambda$ par $J/R\longrightarrow G$.

D'ailleurs, la formation du quotient ainsi \unsure{obtenu|obtenus} $J/R$ de $J$ commute à l'extension plate de la base $S$ [en particulier à la localisation].

Ceci assure la validité des conditions \uline{7$^\circ$} et \uline{8$^\circ$}.

\noindent\rule{\linewidth}{0.3pt}

\uline{Corollaires du \unsure{th|Th}}

\uline{Corollaire 1} Soient $P/X/k$ sur $k$ alg. clos \inserted{(\uline{sép. clos suffit}!)} $X$ propre sur $k$, $P/X$ loc de \deleted{\unsure{présentation}}\inserted{type} fini. Alors $\prod_{X/k}P/\deleted{X}$ représentable par préschéma loc \deleted{\unsure{de}} algébrique.

Si $P/X$ quasi-projectif\inserted{ ou seulmt à fibres \unsure{qs. pro.|qu.-pr.}}, alors $\prod_{X/k}\longrightarrow{}P/k$ est essentiellement quasi-projectif [i.e. tt ouvert quasi-compact est quasi-projectif] sur $k$.

\unsure{$\swarrow$} il suffit que pour tt $x\in\operatorname{Ass}\mathcal O_X$, $\exists\,y\in\bar x$ tel que la fibre $P_y$ soit quasi-projective.

% archive sheet number: 24
